% Copyright (c) 2026 Eugene Joseph M. Ragasa.
% SPDX-License-Identifier: Apache-2.0
% Author: Eugene Joseph M. Ragasa
% Formal proof completed: 2026-10-31
\documentclass[11pt]{article}
\usepackage[T1]{fontenc}\usepackage{amsmath,amssymb}\usepackage[margin=1in]{geometry}\usepackage[hidelinks]{hyperref}
\title{Proof Dossier: Finite Matrix-Product Convolution}\author{Eugene Joseph M. Ragasa}\date{Formal proof completed: 2026-10-31}
\begin{document}\maketitle
\section{Identity and status}
ID \texttt{BG-FOURIER-FINITE-THM-001}; status \texttt{proved}. Lean declaration:
\par\noindent{\tiny\path{ProjectKoios.BandGauge.Fourier.FiniteAbelian.finiteFourierCoefficient_mul}}.
\section{Exact statement and proof}
For a finite abelian group and compatible finite complex matrix fields $L$ and $R$, the
normalized conjugated-character Fourier coefficient of the pointwise product $L(k)R(k)$ equals
\[
  (\widehat L\star\widehat R)(m)
  =\sum_{p,q}\mathbf 1_{p+q=m}\,\widehat L(p)\widehat R(q).
\]
Reconstructing the convolution distributes the character phase over both finite sums. Character
multiplicativity turns the product of phases into the phase at $p+q$, and the conditional sum
selects that mode. Fourier inversion then identifies the coefficient of the pointwise product.
\section{Conventions and boundary}
The forward transform carries $|G|^{-1}$ and a conjugated character; reconstruction uses the
unconjugated character and no additional normalization. Consequently no extra cardinality factor
appears. The result is exact finite algebra and does not bind an FFT, runtime array, external mesh,
continuum transform, or physical hopping range.
\section{Trusted base and verification}
Checked with Lean \texttt{4.35.0} and Mathlib revision
\texttt{56a005855ea17e0b0958f73ec801ae140b0a1679}. Kernel axioms are
\texttt{propext}, \texttt{Classical.choice}, and \texttt{Quot.sound}; no admissions.
AI-assisted and kernel-checked; no independent expert review recorded.
\end{document}
